\section{Introduction}
\label{sec:intro}

Modern biological research depends on precise, repeatable laboratory protocols --- yet the systems that execute them are largely incapable of self-correction.
When a pipetting robot sets the wrong CO$_2$ concentration, or a reagent is prepared at an incorrect dilution, current automation systems either fail silently or halt and wait for human intervention.
This brittleness is a fundamental bottleneck in high-throughput biology: a single protocol error can invalidate entire experimental batches, wasting reagents, time, and samples.
An AI agent capable of detecting and correcting such errors in real time would directly address this bottleneck, enabling continuous autonomous experimentation.

{\bf Why is protocol error correction hard?}
Biological protocols combine numerical parameters (concentrations, temperatures, speeds), procedural logic (step ordering, timing), and domain-specific terminology that interacts in subtle ways.
Correctly fixing a protocol error requires not only identifying which parameter is wrong, but also knowing the biologically appropriate replacement value --- and making \emph{only that change} while leaving the rest of the protocol intact.
Recent large language models (LLMs) have shown strong reasoning capabilities in scientific domains~\cite{bran2023chemcrow,schmidgall2025agentlab}, but their application to wet-lab error correction remains largely unexplored.

{\bf Gap in existing work.}
Prior systems for autonomous scientific experimentation fall into two categories.
Chemistry-focused tool-augmented agents (\eg \chemcrow~\cite{bran2023chemcrow}) demonstrate that LLMs with domain tools can plan and execute synthesis protocols, but they operate on single-step reaction planning rather than multi-parameter protocol validation.
Autonomous research pipeline agents (\eg \agentlab~\cite{schmidgall2025agentlab}) address the full research lifecycle but focus on computational machine learning experiments rather than physical wet-lab execution.
In the biology domain, \bioprobench~\cite{bioprobench2025} provides a comprehensive benchmark for protocol understanding, but no prior work has systematically evaluated AI agents on the error correction subtask with closed-loop feedback adaptation.
Crucially, the failure modes of LLMs on protocol correction --- particularly the tendency to over-modify protocols --- have not been characterized.

{\bf Our approach.}
We propose \wetlabagent, a modular AI agent with four experimental configurations of increasing sophistication:
(1) a \vanilla baseline that issues a single-shot correction prompt,
(2) a \toolaug agent that runs domain-specific parameter validators before querying the LLM,
(3) a \feedbackone agent that performs one round of simulated execution feedback, and
(4) a \feedbackthree agent that performs three rounds of closed-loop adaptation.
We evaluate all four conditions on the \bioprobench error correction task using \claude, measuring error detection accuracy, exact-match correction accuracy, and semantic correction accuracy.

{\bf Key findings.}
All conditions achieve 80--90\% error detection accuracy, demonstrating that LLMs are reliable protocol error \emph{detectors}.
However, exact-match correction accuracy is only 20--30\% across conditions.
We show that this gap arises from a systematic over-correction failure mode: the LLM correctly identifies and fixes the primary error but also reformats surrounding text, causing exact-match failures even when the scientific fix is correct.
Semantic accuracy --- which measures whether the primary fix is correct regardless of formatting --- reaches 70--80\%, with \feedbackthree achieving the highest value of 80\%.
Operation-type errors (\eg incorrect centrifuge speeds) are the hardest, with only 50\% semantic accuracy even for the best agent.

Our contributions are:
\begin{itemize}[leftmargin=*,itemsep=0pt,topsep=0pt]
    \item We propose \wetlabagent, a modular agent framework for biological protocol error detection and correction, evaluating four conditions from single-shot to three-round closed-loop adaptation.
    \item We identify and characterize the \emph{over-correction} failure mode in LLM-based protocol correction --- a previously underreported phenomenon with direct implications for wet-lab deployment safety.
    \item We conduct a systematic per-error-type analysis on \bioprobench, revealing that operation errors (centrifuge speed, timing) are significantly harder than parameter or reagent errors.
    \item We demonstrate that closed-loop feedback adaptation provides measurable semantic accuracy gains (70\% $\to$ 80\%) and recommend minimal-change prompt engineering as a practical mitigation for over-correction.
\end{itemize}
